\section{Conclusion}
\label{sec:conclusion}

This study presented MindTrace, a unified benchmarking framework for multi-class emotion classification that systematically compared two traditional machine learning models (SVM, XGBoost) and two deep learning architectures (CNN, BiLSTM) on a 416,809-sample Twitter corpus under unified preprocessing, stratified splits, and macro-averaged F1 evaluation. Regarding \textbf{RQ1}, BiLSTM achieved the highest test accuracy (94.70\%, macro-F1 0.9464), followed by CNN (93.13\%), SVM (91.80\%), and XGBoost (80.93\%), suggesting that bidirectional contextual encoding provides a measurable advantage over local feature extraction and bag-of-words representations, with XGBoost's substantially lower accuracy indicating that gradient-boosted trees over TF-IDF features struggle with fine-grained emotion discrimination. For \textbf{RQ2}, confusion matrix analysis identified Joy$\leftrightarrow$Love (valence conflation) and Fear$\rightarrow$Surprise (arousal ambiguity) as architecture-independent misclassification patterns, establishing that semantic overlap --- rather than model capacity --- defines the current classification boundary. Addressing \textbf{RQ3}, SVM was selected over BiLSTM due to its ${\sim}23\times$ lower inference latency and was integrated into a production-ready Flask REST API with Docker containerisation and CI/CD automation, demonstrating that a classical model with TF-IDF features can deliver reliable real-time emotion prediction at minimal computational cost; ablation studies further confirmed that stopword removal ($-1.33$\,pp) and BiLSTM depth ($-1.22$\,pp) are the most impactful pipeline components. The complete codebase is publicly available to support reproducibility.

Several limitations warrant acknowledgement. The evaluation relies on a single stratified train--test split without $k$-fold cross-validation, and no statistical significance testing (\eg, paired bootstrap or McNemar's test) was conducted; consequently, the observed performance differences may reflect split-specific variance rather than true model superiority. Class imbalance was addressed through random downsampling, potentially discarding informative samples; alternative strategies such as class weighting or oversampling merit investigation. The dataset is drawn exclusively from Twitter, and both platform-specific language patterns and crowd-sourced annotation noise may limit generalisability to longer-form domains. Hyperparameter tuning was intentionally constrained to ensure fair comparison, and transformer-based architectures (BERT, RoBERTa) were excluded due to computational constraints, leaving the comparison incomplete with respect to the current state of the art.

Future work should adopt cross-validation and significance testing, integrate transformer-based models alongside additional baselines (\eg, logistic regression, Na\"ive Bayes), and explore cross-domain transfer to non-Twitter corpora. Enhanced preprocessing --- such as emoji-to-text conversion, punctuation-based intensity modelling, and elongated-word normalisation --- along with interpretability methods (SHAP, LIME), qualitative error analysis, and ethical safeguards around data anonymisation and fairness represent important directions toward a more comprehensive emotion classification framework.
